% Nag is used for debug
\RequirePackage[l2tabu, orthodox]{nag}
\documentclass[a4paper,french,twoside,13pt]{report}
%\documentclass[a4paper,twoside,10pt]{article}
%\usepackage[bookman]{quotchap}

\usepackage[bindingoffset=1.5cm,centering,includeheadfoot,margin=3cm]{geometry}

%\frenchspacing
\usepackage[frenchb]{babel}
\usepackage[T1]{fontenc}
\usepackage[utf8]{inputenc}
\usepackage[pdftex]{graphicx}
\usepackage{subfig}

\usepackage[protrusion=true,expansion=true]{microtype}

\usepackage{fancybox}
\usepackage{shadow}
\usepackage{xspace}
\usepackage{acronym}

% Maths enhancing
\usepackage{amsmath}
\usepackage{amssymb}
\usepackage{amsthm}
\usepackage{dsfont}

\usepackage{algorithm2e}

% JP : automatic line return after paragraph
\makeatletter
\renewcommand\paragraph{\@startsection{paragraph}{4}{\z@}%
    {-3.25ex\@plus -1ex \@minus -.2ex}%
    {1.5ex \@plus .2ex}%
{\normalfont\normalsize\bfseries}}
\makeatother

% JP : use amsthm package and define some useful commands
\newtheorem{theorem}{Theorem}
\newtheorem{definition}{Definition}
\newtheorem{proposition}{Proposition}

%\usepackage[right]{eurosym}
\usepackage{lmodern} % new latin extended computer modern font}, use with T1
%\usepackage{pxfonts} % has replaced palatino
%\usepackage{cmbright}

\usepackage{listings}

%\usepackage{booktabs}
%\usepackage{array}

\usepackage[allowlitunits]{siunitx}
% Try $20\micro\seconds$

%\addtolength{\oddsidemargin}{-1cm}
%\addtolength{\evensidemargin}{-1cm}
%\addtolength{\textwidth}{+2cm}
%\usepackage{showframe}

%\usepackage{lastpage}
% for \pageref{LastPage}

% fancy header
\usepackage{fancyhdr}
\pagestyle{fancy}
%\pagestyle{plain} % Imprime simplement le numero de page centre en bas
%\pagestyle{headings}
\fancyhf{} % clear default header and footer

\fancyhead[LE,RO]{\textit{\nouppercase{\leftmark}}}
%\fancyhead[LO,RE]{\textit{\nouppercase{\thepage{} of \pageref{LastPage}}}}
\fancyhead[CE,CO]{}
\fancyfoot[C]{--\thepage--}
\setlength{\headheight}{14pt}

\newcommand\Hrule{\noindent\rule{\linewidth}{1.5pt}}
%\AL: to add abs and norm for fracs
%\DeclarePairedDelimiter\abs{\lvert}{\rvert}%
%\DeclarePairedDelimiter\norm{\lVert}{\rVert}%
\usepackage[pdftex,
    pdfauthor={Alexandre Lafaye},
    pdftitle={Laboratoire PHY1 E7},
    pdfsubject={},
    pdfkeywords={rapport,laboratoire,HEIG-VD,PYH1,Thérmodynamique},
pdfcreator=pdflatex]{hyperref}
\hypersetup{colorlinks, citecolor=black, filecolor=black, linkcolor=black, urlcolor=black}

\newcommand{\weblink}[2]{\href{#1}{\url{#1}: \textsc{#2}}}

%Maths shortcuts
%\newcommand{\R}{\mathbb{R}}
%\newcommand{\E}{\mathbb{E}}
%\newcommand{\N}{\mathbb{N}}
%\newcommand{\Var}{\mathrm{Var}}
%\newcommand{\Cov}{\mathrm{Cov}}
%\newcommand{\Prob}{\mathbb{P}}
%\newcommand{\trace}{\mathbb{\text{trace}}}
%\def\mathbi#1{\textbf{\em #1}}
%%\def\mathi#1{\em #1}
%\newcommand{\Z}{\mathbi{Z}_2}
%\newcommand{\kMeans}{{\em k}-Means\xspace}
%\newcommand{\kmeans}{{\em k}-means\xspace}
%\newcommand{\kmeanspp}{{\em k}-means$++$\xspace}
%\newcommand{\captionwidth}{0.9\textwidth}

% index stuff
\usepackage{tocbibind} % Index in TOC
%\usepackage{makeidx}
%\usepackage{showidx} % show index entries in margin
%\makeindex

%\newcommand{\bib}[4]{\item{\textsc{#1}, \emph{#2}, #3\\\hspace{2em}#4}}

% ---------------------------------------------------------

\newcommand{\titleinfo}{Laboratoire de Physique E7 }
\author{Alexandre Lafaye}
\date{07 Avril 2016}
\title{\titleinfo}

% Margins for handwritten notes
%\addtolength{\oddsidemargin}{-2cm}
%\addtolength{\evensidemargin}{-2cm}
%\addtolength{\textwidth}{-2cm}

\usepackage{color}
\definecolor{Brown}{cmyk}{0,0.81,1,0.60}
\definecolor{OliveGreen}{cmyk}{0.64,0,0.95,0.40}
\definecolor{CadetBlue}{cmyk}{0.62,0.57,0.23,0}
\definecolor{lightlightgray}{cmyk}{0,0,0,0.02}
\definecolor{gray}{cmyk}{0,0,0,0.8}

% LaTeX magic: make sections and parts have a cleardoublepage
%\let\stdsection\section
%\renewcommand*\section{\cleardoublepage\stdsection}
\captionsetup{belowskip=12pt,aboveskip=4pt}
\begin{document}
% JP
%\frontmatter

\lstset{
    language=Python,                             % Code langugage
    basicstyle=\ttfamily,                   % Code font, Examples: \footnotesize, \ttfamily
    keywordstyle=\color{OliveGreen},        % Keywords font ('*' = uppercase)
    commentstyle=\color{gray},              % Comments font
    numbers=left,                           % Line nums position
    numberstyle=\tiny,                      % Line-numbers fonts
    stepnumber=1,                           % Step between two line-numbers
    numbersep=5pt,                          % How far are line-numbers from code
    backgroundcolor=\color{lightlightgray}, % Choose background color
    frame=none,                             % A frame around the code
    tabsize=2,                              % Default tab size
    captionpos=b,                           % Caption-position = bottom
    breaklines=true,                        % Automatic line breaking?
    breakatwhitespace=false,                % Automatic breaks only at whitespace?
    showspaces=false,                       % Dont make spaces visible
    showtabs=false,                         % Dont make tabls visible
    morekeywords={uint32_t,uint8_t,uint16_t,time_spec_t,size_t,
    clock_config_t},
}
\pagestyle{empty}
\pagenumbering{arabic}
\begin{titlepage}
    \null\vfil
    \begin{flushleft}
	\begin{center}
	    \Huge \textbf{\titleinfo}
	\end{center}
    \end{flushleft}
    \par
    \hrule height 4pt
    \par
    \begin{flushright}
	\large
	\textsl{Rapport de laboratoire} \par
    \end{flushright}
    \vspace{\fill}
    \vspace*{\stretch{1}}

    \begin{center}
	\Large
	\textsc{Colin Bairoch, Julia Cosandier, Alexandre Lafaye\\
	    %\normalsize
	HEIG-VD -- 2016} \\
	\href{mailto:julia.cosandier@heig-vd.ch}{julia.cosandier@heig-vd.ch}
	\href{mailto:colin.bairoch@heig-vd.ch}{colin.bairoch@heig-vd.ch}
	\href{mailto:alexandre.lafaye@epfl.ch}{alexandre.lafaye@epfl.ch}
    \end{center}
    \vspace{\fill}
    \vspace*{\stretch{2}}
    \begin{center}
	\large
	Supervisé par: \\
	\vspace{\fill}
	Ioan Balin -- HEIG-VD\\
	\href{http://heig-vd.ch/}{http://heig-vd.ch/} \\
    \end{center}
    \begin{center}
	\large
	Mirko Croci -- HEIG-VD \\
	\href{http://heig-vd.ch/}{http://heig-vd.ch/} \\
    \end{center}

    \vspace{\fill}
    \vspace*{\stretch{2}}

    \begin{figure}[!h]
	\centering
	\parbox{1.4in}{
	\includegraphics[height=4em]{logo_heig_h80.jpg}}
    \end{figure}

    \vspace*{1cm}
\end{titlepage}

\pagestyle{fancy}

\tableofcontents

\newpage

\section{Introduction}
%\addcontentsline{toc}{section}{Introduction}
\markboth{Introduction}{}
\subsection*{Objectifs du Laboratoire}
Dans ce travail pratique les buts suivants ont été visés:
\begin{itemize}
    \item{Détermination de la chaleur massique de différentes substances}
	\begin{itemize}
	    \item{Cuivre}
	    \item{Plomb}
	    \item{Verre}
	\end{itemize}
    \item{Détermination de la chaleur latente de fusion et de vaporisation de l'eau}
\end{itemize}


\section{Théorie}
%\addcontentsline{toc}{section}{Théorie}
\markboth{Théorie}{}
\subsection{Chaleur massique}
Lors ce que l'on chauffe ou on refroidit un corps, l'énergie absorbée sous forme de chaleur $\Delta Q$ est proportionnelle au changement de température $\Delta \theta$ et la masse m du corps:
\begin{equation}
    \Delta Q = m\cdot c \cdot \Delta \theta
\end{equation}

Où c représente la chaleur massique, un coefficient dépendant de la nature du matériau. Pour le déterminer, 

Quantité de chaleur cédée par la grenaille:
\begin{equation}
    \label{deltaq}
    \Delta Q_1 = m_1 \cdot c_1 \cdot (\theta_1 -\theta_M)
\end{equation}
qui est égale à celle absorbée par l'eau.
\begin{equation}
    \Delta Q_2 = m_2 \cdot c_2 \cdot (\theta_M -\theta_2)
\end{equation}
On suppose ici qu'on connaît déjà la valeur de la chaleur massique de l'eau $c_2$, et $\theta_1 = \sim 98C\degree$ est égal à la température de la vapeur d'eau.
En utilisant les deux formules précédentes:
\begin{equation}
    c_1 = c_2\cdot \frac{m_2\cdot (\theta_M - \theta_2)}{m_1 \cdot (\theta_1 - \theta_M)}
\end{equation}

Il faut aussi prendre en compte que le vase calorimétrique dans lequel le mélange va être fait. La quantité de chaleur qu'il va absorber peut être donnée en masse équivalente d'eau: $m_{cal}=0.023kg$.
\begin{equation}
    \Delta Q_2 = (m_2+m_{cal}) \cdot c_2 \cdot (\theta_M -\theta_2)
\end{equation}
Ce qui donne donc:
\begin{equation}
    \label{final_cm}
    c_1 = c_2\cdot \frac{(m_2 + m_{cal})\cdot(\theta_M - \theta_2)}{m_1 \cdot (\theta_1 - \theta_M)}
\end{equation}

\subsection{Chaleur latente de vaporisation}
Quand un corps plongé dans un milieu où la pression est fixe est réchauffé, sa température augmente. Sauf si la substance du corps est dans une transition entre deux phases, toute l'énérgie apportée est alors dissipée pour passer d'une phase à l'autre. Dès que le chagement est terminé, la température va de nouveau augmenter.
Une bonne manière d'observer ce phénomène est en 

La chaleur est donc partagée entre la partie absorbée par~\eqref{deltaq}.

\begin{equation}
    \Delta Q_2 = m_1 \cdot L_v
\end{equation}


\begin{equation}
    \Delta Q_3 = m_2 \cdot c_e \cdot (\theta_M - \theta_2)
\end{equation}

\begin{equation}
    \Delta Q_4 = m_{cal} \cdot c_e \cdot (\theta_M - \theta_2)
\end{equation}

$\Delta Q_3 + \Delta Q_4$

\begin{equation}
    m_1 \cdot c_e \cdot (\theta_V - \theta_M) + m_1 \cdot L_V = (m_2 + m_{cal}) \cdot c_e \cdot (\theta \theta)
\end{equation}

\begin{equation}
    \label{final_lv}
    L_V = c_e \cdot \frac{(m_2 + m_{cal})\cdot (\theta_M - \theta_2) - m_1 \cdot (\theta_V - \theta_M)}{m_1}
\end{equation}

\subsection{Chaleur latente de fusion}
\begin{equation}
    \Delta Q_1 = m_1 \cdot c_e \cdot (\theta_M - \theta_1)
\end{equation}

\begin{equation}
    \Delta Q_2=m_1 \cdot L_F
\end{equation}

\begin{equation}
    \Delta Q_3 = m_2 \cdot c_e \cdot (\theta_2 - \theta_M)
\end{equation}

\begin{equation}
    \Delta Q_4 = m_{cal} \cdot c_e \cdot (\theta_2 - \theta_M)
\end{equation}

\begin{equation}
    \label{final_lf}
    L_F = c_e \cdot \frac{(m_2 + m_{cal}) \cdot (\theta_2 -\theta_M)-m_1 \cdot (\theta_M - \theta_1)}{m_1}
\end{equation}

\subsection{Incertitudes}

Une mesure expérimentale est toujours accompagnée de son $\textit{incertitude de mesure}$. On peut qualifier cette incertitude selon différentes caractéristiques:
\begin{itemize}
\item La résolution
\item La précision
\item La reproductibilité
\end{itemize}

Cette incertitude a de multiples sources, humaines ou liées au matériel, qui la rendent inévitable mais pas pour autant non-quantifiable.
On la calcule en générale de deux manières, sous forme $\textit{d'incertitude absolue}$ ou $\textit{d'incertitude relative}$.

\paragraph{Notation}
Un résultat s'écrit donc sous la forme: \[a\pm\Delta a\] tel que $\Delta a$ soit l'incertitude absolue et $\frac{\Delta a}{a} \%$ soit l'incertitude relative.

\paragraph{Chiffres significatifs}
Pour noter correctement les résultats avec leur incertitude, il faut être attentif aux $\textit{chiffres significatifs}$, les mesures et résultats de calculs doivent être exprimés avec un ou x chiffres dont la valeur n'est pas certaine.\\
Par exemple, pour une mesure au gramme près $m = (2.3\pm0.1)kg$.

\paragraph{Propagation des erreurs}
Une fois l'incertitude estimée, il faut encore la propager à tous les calculs qui suivront. Dans le cas d'une fonction à une seule variable, on peut utiliser sa différentielle:
\begin{equation}
    dg=|{\frac{df}{dx}}|dx
\end{equation}

Ce qui donne pratiquement:
\begin{equation}
    \Delta g=|\frac{df}{dx}|\Delta x
\end{equation}

Pour propager les erreurs durant les calculs, les formules suivantes sont utilisées:\\
Addition/Soustraction:
\begin{equation}
    \begin{split}
	g = a \pm b\\
	\Delta g = \Delta a + \Delta b
    \end{split}
\end{equation}
Multiplication/Division:
\begin{equation}
    \label{multdiv}
    \begin{split}
	g = k\cdot ab \;\;\; k=constante\\
	\frac{\Delta g}{|g|} = \frac{\Delta a}{|a|} + \frac{\Delta b}{|b|}
    \end{split}
\end{equation}


\section{Manipulation}
%\addcontentsline{toc}{section}{Manipulation}
\markboth{Manipulation}{}
\subsection{Détermination de la chaleur massique d'un solide}

\begin{figure}[!h]
    \centering
    \caption{Photo Montage chaleur massique}
    \includegraphics[totalheight=26em]{pic/scheme.jpg}
\end{figure}
\begin{enumerate}
    \item{Thermomètre}
    \item{Pipette d'eau distillée}
    \item{Grenaille de plomb}
    \item{Grenaille de verre}
    \item{Grenaille de cuivre}
    \item{Bécher}
    \item{Étuve}
    \item{Tuyau en silicone}
    \item{Vase calorimétrique}
    \item{Générateur de vapeur}
    \item{Balance}
    \item{Piège à eau}
\end{enumerate}

Pour cette partie de l'expérience, il s'agit de faire chauffer trois matériaux différents à une température fixée (celle de la vapeur d'eau), et de les plonger ensuite dans un calorimètre rempli d'eau. La différence de température de l'eau permettra ensuite de savoir quelle quantité de chaleur était stockée dans le matériau en question.\\
\newpage
Les trois matériaux proposés sont:
\begin{enumerate}
    \item{Le cuivre}
    \item{Le plomb}
    \item{Le verre}
\end{enumerate}
De l'eau est placée dans un récipient fermable de manière étanche, ce récipient est ensuite relié à une étuve à l'aide d'un tuyeau en silicone. L'étuve est elle-même branchée sur un autre récipient permettant de récuprer l'eau de la condensation.\\
Après avoir été soigneusement pesé, une certaine quantité du matériau est déposé dans l'étuve grâce à un panier prévu à cet effet. Le récipient contenant l'eau est alors chauffé jusqu'à ce que toute l'eau se soit évpaporée. Le matériau devrait être aux alentours de 98C\degree.\\
Durant la pérdiode de chauffe, le calorimètre a été rempli avec 100g d'eau désionisée à température ambiante.
Finalement, le matériau a été versé dans le calorimètre, ce dernier refermé hérmétiquement le temps que la température interne arrive à son point d'équilibre.
La température finale du mélange entre eau et matériau a été relevée.

Pour calculer la chaleur massique dans chaque cas, la formule~\eqref{final_cm} a été utilisée.
Quand aux incertitudes, la méthode des dérivées partielles s'est révélée utile pour les calculer:

\begin{align*}
    \Delta c_1 = |\frac{\delta c_1}{\delta m_1}|\cdot \Delta m_1 + |\frac{\delta c_1}{\delta m_2}|\cdot \Delta m_2 + |\frac{\delta c_1}{\delta \theta_m}|\cdot \Delta \theta_m + |\frac{\delta c_1theta_2}{\delta \theta_2}|\cdot \Delta \theta_2
\end{align*}

Donc:

\begin{align*}
    \Delta c_1 &= c_2 \frac{(m_2 + m_{cal})\cdot (\theta_m - \theta_2)}{m_1^2 \cdot (\theta_1 - \theta_m)}\cdot \Delta m_1\\
    &+ c_2 \frac{(\theta_m - \theta_2)}{m_1\cdot (\theta_1 - \theta_m)}\cdot \Delta m_2\\
    &+ c_2 \frac{(m_2 + m_{cal})\cdot ((\theta_1 - \theta_m)+(\theta_m - \theta_2))}{m_1 \cdot (\theta_1 - \theta_m)^2} \Delta \theta_m\\
    &+ c_2 \frac{(m_2 + m_{cal})}{m_1 \cdot (\theta_1 - \theta_m)}\Delta \theta_2
\end{align*}

Avec nos mesures, les résultats suivants ont été obtenus:

\begin{table}[!h]
    \centering
    \caption{Résultats Cuivre}
    \begin{tabular}{|l|l|l|}
	\hline
	Mesure	&Valeur	&$\Delta$\\
	\hline
	$m_1$ [kg]	&0.0627	&10e-5\\
	$m_2$ [kg]	&0.1	&10e-5\\
	$m_{cal}$ [kg]	&0.023	&0\\
	$\theta_1$ [C\degree]	&98	&0\\
	$\theta_2$ [C\degree]	&20.8	&0.2\\
	$\theta_m$ [C\degree]	&23.8	&0.2\\
	$c_2$ [J/ (kg*K)]	&4180	&0\\
	\hline
    \end{tabular}
\end{table}

La valeur de la chaleur massique de l'eau ainsi que celles de la masse du calorimètre $m_{cal}$ et celle de la température de la vapeur d'eau ont été considérées comme étant parfaites, car ne provenant pas de nous, mais de tables ou du manuel d'expérience.

Avec la formule~\eqref{final_cm} et celles présentées ci-dessus pour les incertitudes, la valeur suivante a été obtenue pour la chaleur massique du cuivre:
\begin{equation}
    c_1=331.53\pm 45.18[J/ (kg*K)]
\end{equation}
L'incertitude sur ce résultats est assez grande (plus de 13\%) mais en comparant avec la valeur de la table CRM, l'observation est plutôt que cette dernière est relativement éloignée de notre résultat:\\
\begin{equation}
    Valeur\quad tabulée:\quad c_1 = 390[J/ (kg*K)]
\end{equation}

Les grosses sources d'erreurs sur cette manipulation viennent probablement de la fuite de chaleur que représentent les tuyeau de silicone. En effet, une partie de l'eau est restée condensée dans les tuyeaux et leur température a grandement augmenté durant l'expérience.
Une autre source est que le cuivre a du refroidir au contact du panier et lui communiquer une partie de sa chaleur.
Le calorimétre devant laisser entrer le thermomètre, il n'est pas étanche, et un phenomène de convection de l'air à l'intérieur de ce dernier est probablement une source de fuite d'énérgie.
Il est aussi impossible de savoir à quelle température se trouvait exactement le cuivre en sortant de l'étuve, la valeur de 98C\degree n'étant qu'une supposition raisonnable.\\
Pour toutes les mesures qui vont suivre, $m_{cal}$ est la masse équivalent en eau du calorimétre.

Pour le plomb les mesures ont été:
\begin{table}[!h]
    \centering
    \caption{Résultats Plomb}
    \begin{tabular}{|l|l|l|}
	\hline
	Mesure	&Valeur	&$\Delta$\\
	\hline
	$m_1$ [kg]	&0.0265	&10e-5\\
	$m_2$ [kg]	&0.1001	&10e-5\\
	$m_{cal}$ [kg]	&0.023	&0\\
	$\theta_1$ [C\degree]	&98	&0\\
	$\theta_2$ [C\degree]	&21.2	&0.2\\
	$\theta_m$ [C\degree]	&22.6	&0.2\\
	$c_2$ [J/ (kg*K)]	&4180	&0\\
	\hline
    \end{tabular}
\end{table}

En réutilisant la formule~\eqref{final_cm}:
\begin{equation}
    c_1 = 360.53 \pm 104.13[J/ (kg*K)]
\end{equation}
Mais la valeur tabulée reste relativement loin de notre résultat:
\begin{equation}
    Valeur\quad tabulée:\quad c_1 = 129[J/ (kg*K)]
\end{equation}
En plus des des sources d'incertitudes évoquées lors du cas du cuivre, lors de cette manipulation deux billes de plomb se sont échappées lors du transfert dans le calorimétre, provoquand ainsi probablement la grande différence entre les valeurs `officielles' trouvées dans le CRM et notre mesure.

\newpage
Et finalement pour le verre:
\begin{table}[!h]
    \centering
    \caption{Résultats Verre}
    \begin{tabular}{|l|l|l|}
	\hline
	Mesure	&Valeur	&$\Delta$\\
	\hline
	$m_1$ [kg]	&0.0153	&10e-5\\
	$m_2$ [kg]	&0.1	&10e-5\\
	$m_{cal}$ [kg]	&0.023	&0\\
	$\theta_1$ [C\degree]	&98	&0\\
	$\theta_2$ [C\degree]	&21.2	&0.2\\
	$\theta_m$ [C\degree]	&23.8	&0.2\\
	$c_2$ [J/ (kg*K)]	&4180	&0\\
	\hline
    \end{tabular}
\end{table}

Les mesures ici semblent un peu meilleure, car:
\begin{equation}
    c_1 = 1177.49 \pm 185.19[J/ (kg*K)]
\end{equation}
Et:\\
\begin{equation}
    Valeur\quad tabulée:\quad c_1 = [830,1450][J/ (kg*K)]
\end{equation}
Ne sachant pas de quel type de verre il s'agit exactement, on ne dipose pas de mieux qu'une certaine portée de résultats possibles dans lequel est compris le résultat de ce laboratoire. Les sources d'erreurs sont les même que pour les cas précédents.

\newpage
%%%%%%%%%%%%%%%%%%%%%%%%%%%%%%%%%%%%%%%%%%%%%%%%%%%%%%%%%%%%%%%%%%%%%%%%%%%%%%%%
\subsection{Chaleur latente de vaporisation}
Pour cette partie de la manipulation, la but est d'observer le phenomène de vaporisation et de mesurer la quantité d'énérgie échangée lors de cette transformation afin de pouvoir calculer la valeur de la chaleur latente de vaporisation de l'eau.
Pour ce faire, le montage expérimental suivant a été utilisé:

\begin{figure}[!h]
    \centering
    \caption{Photo montage vaporisation}
    \includegraphics[totalheight=21em]{pic/figure10.JPG}
\end{figure}

Le générateur de vapeur a préalablement été rempli d'eau distillée et soigneusement refermé. Il a été ensuite relié au piège à eau avec un tube en silicone assurant l'étanchéité du système.\\
Dans un premier temps le piège à eau a été placé au dessus d'un bêcher, le temps que le système se stabilise un peu, permettant ainsi d'éviter les éclaboussure de vapeur.\\
Dans un second temps, le piège à eau a été plongé dans le calorimètre rempli de 150g d'eau préalablement pesé. La balance ayant été réglée pour une valeur du poids du calorimètre plus 20g elle arrive à l'équilibre au moment ou cette masse d'eau s'est condensée dans le calorimètre.\\
La température a alors été mesurée.

Une fois de plus les dérivées partielles ont été exploitées pour connaître les incertitudes ainsi que la formule~\eqref{final_lv}:
\begin{align*}
    \Delta L_v = |\frac{\delta L_v}{\delta m_1}|\cdot \Delta m_1 + |\frac{\delta L_v}{\delta m_2}|\cdot \Delta m_2 + |\frac{\delta L_v}{\delta \theta_m}|\cdot \Delta \theta_m + |\frac{\delta L_vtheta_2}{\delta \theta_2}|\cdot \Delta \theta_2
\end{align*}

Donc:

\begin{align*}
    \Delta L_v &= c_e \frac{(m_2 + m_{cal})\cdot (\theta_m - \theta_2)}{m_1^2}\cdot \Delta m_1\\
    &+ c_e \frac{(\theta_m - \theta_2)}{m_1}\cdot \Delta m_2\\
    &+ c_e \frac{(m_2 + m_{cal})}{m_1} \Delta \theta_m\\
    &+ c_e (\frac{(m_2 + m_{cal})}{m_1} + 1)\Delta \theta_2\\
\end{align*}

Les mesures suivantes ont été relevées:

\begin{table}[!h]
    \centering
    \caption{Résultats Chaleur latente vaporisation}
    \begin{tabular}{|l|l|l|}
	\hline
	Mesure	&Valeur	&$\Delta$\\
	\hline
	$m_1$ [kg]	&0.02	&10e-4\\
	$m_2$ [kg]	&0.15	&10e-5\\
	$m_{cal}$ [kg]	&0.023	&0\\
	$c_e$ [J/ (kg*K)]	&4180	&0\\
	$\theta_2$ [C\degree]	&23.6	&0.2\\
	$\theta_m$ [C\degree]	&82.5	&0.2\\
	$\theta_v$ [C\degree]	&98	&0\\
	\hline
    \end{tabular}
\end{table}

En effectuant les calculs, on trouve: $L_v = 2'064'857.30\pm 121'904.27[J/ (kg*K)]$,
la valeur tabulée étant: $L_v = 2'300'000[J/ (kg*K)]$
Malgré la taille de l'intervalle d'incertitude, les deux valeurs ne se recoupent pas. Il existe toujours pour cette manipulation le problème de la prise de température dans le calorimétre qui permet à la chaleur de s'échapper. Il est aussi possible que le temps nécessaire à attendre pour éviter les éclaboussures de vapeur ne se soit pas écoulé. La mesure devant être prise rapidement 

%%%%%%%%%%%%%%%%%%%%%%%%%%%%%%%%%%%%%%%%%%%%%%%%%%%%%%%%%%%%%%%%%%%%%%%%%%%%%%%%
\subsection{Chaleur latente de fusion}
Cette manipulation est analogue à celle sur la vaporisation dans son principe, mais il s'agit ici d'observer un autre changement d'état qui est la fusion.
Dans ce but, le vase calorimétrique a été rempli de 200g d'eau chaude à environ 50C\degree. Environ 50g de glaçon ont été ajoutés et la mesure de la température du mélange a été efféctuée une fois la glace fondue.

Une fois de plus pour les incertitudes.
En utilisant la formule~\eqref{final_lf}:
\begin{align*}
    \Delta L_f = |\frac{\delta L_f}{\delta m_1}|\cdot \Delta m_1 + |\frac{\delta L_f}{\delta m_2}|\cdot \Delta m_2 + |\frac{\delta L_f}{\delta \theta_m}|\cdot \Delta \theta_m + |\frac{\delta L_ftheta_2}{\delta \theta_2}|\cdot \Delta \theta_2
\end{align*}

Donc:

\begin{align*}
    \Delta L_f &= c_e \frac{(m_2 + m_{cal})\cdot (\theta_2 - \theta_m)}{m_1^2}\cdot \Delta m_1\\
    &+ c_e \frac{(\theta_2 - \theta_m)}{m_1}\cdot \Delta m_2\\
    &+ c_e \frac{(m_2 + m_{cal})}{m_1} \Delta \theta_m\\
    &+ |c_e (\frac{-(m_2 + m_{cal})}{m_1} -1)|\Delta \theta_2\\
\end{align*}

Les mesures obtenues:

\begin{table}[!h]
    \centering
    \caption{Résultats Chaleur latente fusion}
    \begin{tabular}{|l|l|l|}
	\hline
	Mesure	&Valeur	&$\Delta$\\
	\hline
	$m_1$ [kg]	&0.0539	&10e-5\\
	$m_2$ [kg]	&0.1986	&10e-5\\
	$m_{cal}$ [kg]	&0.023	&0\\
	$c_e$ [J/ (kg*K)]	&4180	&0\\
	$\theta_1$ [C\degree]	&0	&0\\
	$\theta_2$ [C\degree]	&55.8	&0.2\\
	$\theta_m$ [C\degree]	&32.9	&0.2\\
	\hline
    \end{tabular}
\end{table}

Ici aussi, la chaleur latente trouvée est inférieure à celle des tables:
\begin{align*}
    &L_f = 256'021.51 \pm 7800.89[J/ (kg*K)]\\
    &Valeur\quad tabulée:\quad L_f = 330'000[J/ (kg*K)]
\end{align*}

Ici aussi, il a peut être été attendu un peu trop longtemps avant de faire le relevé de la mesure, afin d'être certain que l'entier de la glace ait fondu, provoquant une légère hausse de température de l'eau restée au contact avec l'air ambiant.


\section{Conclusion}
%\addcontentsline{toc}{section}{Conclusion}
\markboth{Conclusion}{}
On a pu observer dans ce laboratoire que les résultats obtenus se sont révélés relativement imprécis et les incertitudes mal estimées.


%Not to have clear page...
\listoffigures
\begingroup
\let\clearpage\relax
\listoftables
\endgroup

\end{document}
